\chapter{The report}
\label{ch:report}

Several functions of \fbk{} write a report: \menu{1} writes one for an output
file; \menu{2}, \menu{10} and \menu{11} write one for a structure or a data
file. All reports share the same anatomy, and each comes in two identical-content
flavours: \file{name.fbk.md} for reading and \file{name.fbk.json} for scripts.
The following sections walk through the report of one real run (the example of
\menu{1}, \file{examples/expected/m1\_checkup.txt}).

\section{Summary and analysis sections}

\lstinputlisting[style=fbkcode, firstline=1, lastline=23,
  title={Report head: summary and the first analysis section}]
  {examples/expected/products/n2_casscf_nevpt2.out.fbk.md}

The summary names the subject (your file) and the finding counts by severity.
Then follow the analysis sections, one per analyser that applies to this output.
Here the A1 section assigns every orbital to the (atom, shell) that carries its
largest weight, partitioned by occupation first -- never rank the whole space by
weight, or virtual orbitals will crowd out occupied ones. An A1 line such as
``no f-composition orbital in the occupied partition'' is itself a result: on
an f system it usually means the calculation found a different solution branch
(wrong initial guess or window).

\lstinputlisting[style=fbkcode, firstline=24, lastline=49,
  title={The A2 section: occupation spectrum and the single-orbital entropy bound}]
  {examples/expected/products/n2_casscf_nevpt2.out.fbk.md}

The A2 section prints the occupation spectrum and an entropy \emph{bound}. The
distinction matters (Chapter~\ref{ch:menu1}, and the full derivation in
Chapter~\ref{ch:criteria}): the literature single-orbital entropy of the
autoCAS protocol is defined on the four occupation states of an orbital, and an
ORCA output does not determine it. What the program computes instead is the
maximum-entropy completion of those four weights for the spin-summed
occupation. That bound is never below the true value, so it is used one-sidedly:
a bound at or below 0.14 \emph{excludes} strong multireference character on that
metric; a bound above it decides nothing, and the section says so.

\lstinputlisting[style=fbkcode, firstline=41, lastline=63,
  title={The A3 section: the multi-reference character panel}]
  {examples/expected/products/n2_casscf_nevpt2.out.fbk.md}

The A3 panel collects the available indicators (here: fractional occupations
and the entropy bound; T1 needs a coupled-cluster output), states which of them
fired, which were indecisive and which are missing, and gives one combined
verdict drawn from the available evidence only.

\section{Findings}

Findings are the diagnostics: conclusions from the rule base, ordered by
severity. Every finding names the rule that produced it, the action to take,
the refusals that apply, and its evidence.

\lstinputlisting[style=fbkcode, firstline=68, lastline=92,
  title={Findings of the example run}]
  {examples/expected/products/n2_casscf_nevpt2.out.fbk.md}

The four severities mean:

\begin{center}
\begin{tabularx}{\textwidth}{@{}lX@{}}
\toprule
Severity & Meaning \\
\midrule
\sev{refuse} & the program refuses to draw a conclusion (data missing or contradictory) \\
\sev{error} & a conclusion was reached but must not be relied on \\
\sev{warn} & a concrete doubt was found; act on it before using the numbers \\
\sev{info} & a check that must be carried out (not a detected problem) \\
\bottomrule
\end{tabularx}
\end{center}

The first finding of the example is a \sev{warn} about CASSCF convergence: the
job reported convergence by the energy criterion, which does not imply that the
\emph{orbitals} have converged -- and the perturbation layer depends on the
orbitals, not on the energy. The measured record behind that warning is
carried in the evidence lines: a flat energy with still-drifting orbitals can
move a NEVPT2 result by several kcal/mol.

\section{Provenance and references}

The provenance block lists every source used anywhere in the report, and the
references block prints, for every literature source, the complete citation and
a paste-ready BibTeX entry:

\lstinputlisting[style=fbkcode, firstline=96, lastline=125,
  title={Provenance and the head of the references block}]
  {examples/expected/products/n2_casscf_nevpt2.out.fbk.md}

The reference database is a BibTeX file shipped with the program
(\file{sources.bib}); literature evidence that cannot be resolved in it is
rejected when the rules load. A report is therefore self-contained: every number
in it can be traced to a manual section, a DOI, or a recorded measurement --
including the internal records of this project, which are cited by path where
they exist.

\section{Conventions}

\begin{itemize}
  \item Reports contain no timestamps and no environment-dependent content other
    than the subject path; runs are reproducible.
  \item Sequence fields in the JSON companion (\file{*.fbk.json}) are arrays;
    the JSON is the stable interface for scripting, the Markdown is for humans.
  \item Paths printed inside a report are those of \emph{your} machine
    (Windows-style in this manual's examples); the analysis does not depend on
    the path.
\end{itemize}
